% !TEX root = ../main.tex
\chapter{\cmm{}: recursos avançados}
\label{cap:cmm-avancado}

\section{Interrupção: \texttt{\#PRACA}}

O processador SAPHO tem um pino de interrupção (\code{itr}). Quando ele pulsa, o contador de programa desvia para o ponto do código marcado pela diretiva \code{\#PRACA} --- que, diferentemente das diretivas de configuração, aparece \textbf{no meio do programa}, marcando a linha de retorno. O uso típico é reiniciar o laço de processamento quando o hardware externo sinaliza um novo evento (um \emph{trigger} de aquisição, por exemplo), sem esperar o programa completar a volta.

\section{Sincronização com o hardware: \texttt{\#TOAQUI}}
\label{sec:toaqui}

A diretiva \code{\#TOAQUI}, colocada num ponto do programa, faz o pino \code{cheguei} do processador pulsar sempre que a execução passa por ali. É um farol para o mundo externo: \enquote{cheguei neste ponto do programa}. Serve para sincronizar blocos de hardware com fases do algoritmo --- e é também o mecanismo que a AURORA usa internamente: no \textbf{teste de hardware} (terminal THTEST) e nas simulações com Verilator, a IDE insere um \code{\#TOAQUI} ao final do \code{main()} para detectar o término do programa e encerrar a simulação no instante certo.

\section{FFT e o endereçamento bit-reverso}

A diretiva \code{\#FFTSIZ n} configura o circuito de reversão de bits para FFTs de $2^n$ pontos, usado pela indexação \code{x[k)} (Capítulo~\ref{cap:cmm-tipos-operadores}). Na FFT radix-2, os resultados saem em ordem bit-reversa; com \code{x[k)} a reordenação é gratuita --- o hardware inverte os bits do índice na leitura/escrita, sem instruções extras. O repositório do YANC traz um projeto de exemplo completo (\code{proc\_fft}) com FFT em \cmm{}.

\section{Quanto custa cada construção}
\label{sec:custo-hardware}

O lema \enquote{pague pelo que usar} tem uma consequência prática: \textbf{o custo em FPGA do seu processador é proporcional ao subconjunto da linguagem que você usa}. Durante a compilação, o terminal TASM informa cada bloco de hardware instanciado à medida que os opcodes aparecem pela primeira vez:

\begin{itemize}
  \item chamar uma função (\code{CAL}) instancia a memória de pilha de sub-rotinas;
  \item usar arrays instancia o circuito de endereçamento indexado; \code{x[k)}, o de reversão de bits;
  \item \code{/} instancia o divisor inteiro; \code{\%}, o módulo; as versões \code{float}, os blocos de ponto flutuante correspondentes;
  \item ao final, o TASM resume o percentual da ISA e da ULA efetivamente usados pelo programa.
\end{itemize}

Regras de bolso para hardware enxuto: prefira deslocamentos a divisões por potências de 2 (ou \code{norm()}); prefira \code{mod2()} a \code{abs()} quando só compara magnitudes; evite \code{float} se a dinâmica do sinal cabe em inteiros; e lembre que a primeira ocorrência paga o bloco --- a segunda ocorrência do mesmo operador é \enquote{grátis}.

\section{O caminho C++}
\label{sec:caminho-cpp}

Além da \cmm{}, o YANC compila \textbf{C++} para o mesmo processador, por uma toolchain paralela: o pré-processador \tool{cpppp} (com \code{\#include}, \code{\#define} completo e compilação condicional) seguido do compilador \tool{cppcomp}, que aceita um subconjunto de C99 com extensões de C++ e emite o mesmo assembly intermediário.

O que o caminho C++ oferece que a \cmm{} não tem:

\begin{itemize}
  \item \textbf{ponteiros} (para indexar arrays de tamanho fixo; sem alocação dinâmica);
  \item \textbf{recursão}, com quadros de pilha de verdade;
  \item \textbf{structs} simples (POD) e até recursos de C++ como \emph{lambdas} e destrutores (RAII);
  \item cabeçalhos padrão adaptados: \file{array}, \file{cmath}, \file{cstdint}, \file{vector} e outros.
\end{itemize}

O preço: variáveis locais deixam de ter endereço fixo, e com isso \textbf{deixam de aparecer por nome nas formas de onda} --- o vínculo íntimo entre código e onda, marca do fluxo \cmm{}, se perde em parte. Os parâmetros do processador no fluxo C++ têm padrões próprios (palavra de 32 bits, mantissa 23, expoente 8), ajustáveis por \code{\#pragma yanc}.

\begin{nota}
Recomendação prática: comece todo projeto em \cmm{}. Migre para C++ apenas quando precisar especificamente de recursão, ponteiros ou structs --- e mantenha os processadores de sinal críticos em \cmm{}, onde a visibilidade de simulação é total.
\end{nota}

\section{As mensagens do compilador}
\label{sec:mensagens-compilador}

Os diagnósticos do YANC são bilíngues --- acompanham o idioma da IDE (Capítulo~\ref{cap:configuracoes}) --- e têm personalidade: um erro de sintaxe rende \enquote{Erro de sintaxe na linha N} com um comentário bem-humorado, e esquecer o \code{main()} provoca a pergunta clássica sobre onde ele está. Por trás do humor, a informação é séria; as classes principais:

\begin{longtable}{@{}lp{0.60\linewidth}@{}}
\toprule
\textbf{Classe} & \textbf{Exemplos} \\
\midrule
\endhead
Estrutura & falta de \code{main()}; \code{break}/\code{continue} fora de laço; contagem errada de parâmetros. \\
Nomes & variável inexistente; variável já declarada; uso do \code{i} reservado. \\
Tipos & \code{\%}/\code{norm()} com não-inteiro; intrínsecos não implementados para \code{comp}; \code{++} em complexo. \\
Faixas & estouro de inteiro para o \code{\#NUBITS} escolhido; \code{float} fora da faixa do formato configurado; porta de E/S inexistente. \\
Avisos & conversões implícitas entre tipos (atribuição, parâmetro, retorno); \code{float}/\code{comp} em condição ou índice. \\
\bottomrule
\end{longtable}

Erros aparecem no terminal TCMM com a linha clicável --- o clique leva o editor à linha exata (Capítulo~\ref{cap:terminais}). A Aurora Intelligence também sabe explicar qualquer mensagem do compilador (ferramenta \code{lookup\_compiler\_message}, Capítulo~\ref{cap:aurora-intelligence}).
